\documentclass[11pt,a4paper]{article}
\usepackage[margin=22mm,headheight=15pt]{geometry}
\usepackage{fontspec}
\usepackage{polyglossia}
\setdefaultlanguage{russian}
\setotherlanguage{english}
\setmainfont{cmunrm.otf}
\setsansfont{cmunss.otf}
\setmonofont{cmuntt.otf}
\newfontfamily\cyrillicfont{cmunrm.otf}
\newfontfamily\cyrillicfonttt{cmuntt.otf}
\usepackage{unicode-math}
\setmathfont{latinmodern-math.otf}
\usepackage{xcolor,booktabs,longtable,array,enumitem,graphicx}
\usepackage[most]{tcolorbox}
\usepackage{fancyhdr}
\usepackage[colorlinks=true,linkcolor=blue,urlcolor=blue]{hyperref}
\definecolor{navy}{HTML}{102044}
\definecolor{blue}{HTML}{1168E8}
\definecolor{pale}{HTML}{F3F8FF}
\definecolor{yellow}{HTML}{FFDF55}
\setlength{\parindent}{0pt}
\setlength{\parskip}{5pt}
\setlist[itemize]{leftmargin=*,itemsep=2pt,topsep=4pt}
\setlist[enumerate]{leftmargin=*,itemsep=2pt,topsep=4pt}
\pagestyle{fancy}
\fancyhf{}
\fancyhead[L]{\small\color{navy} Введение в машинное обучение · Пара 1}
\fancyhead[R]{\small\color{blue} Конспект преподавателя}
\fancyfoot[C]{\small\thepage}
\renewcommand{\headrulewidth}{0.4pt}
\newtcolorbox{saybox}[1]{colback=pale,colframe=blue,boxrule=0.6pt,arc=2mm,title={#1},fonttitle=\bfseries,breakable}
\newtcolorbox{warnbox}[1]{colback=yellow!20,colframe=yellow!70!black,boxrule=0.6pt,arc=2mm,title={#1},fonttitle=\bfseries,breakable}
\newcommand{\term}[1]{\textbf{\color{navy}#1}}
\newcommand{\timing}[1]{\hfill{\small\color{blue}#1}}
\newcolumntype{P}[1]{>{\raggedright\arraybackslash}p{#1}}

\begin{document}
\begin{titlepage}
\vspace*{25mm}
{\color{blue}\bfseries\Large ДГТУ · Введение в машинное обучение\par}
\vspace{10mm}
{\color{navy}\bfseries\fontsize{30}{35}\selectfont Первая пара\\[3mm] Курс, оценивание и введение в ML\par}
\vspace{12mm}
{\Large Сценарий преподавателя\par}
\vfill
{\large 90 минут · 31 основной слайд · 8 дополнительных\par}
{\small Версия: 20 октября 2026\par}
\end{titlepage}
\tableofcontents
\newpage
\section{Маршрут первой пары}
План курса перечислен без привязки к неделям. Освободившееся время отведено переобучению, bias–variance и утечкам данных. Завершаем одним слайдом итогов.

Критерии учебных медалей и частичного начисления баллов публикуются до старта соревнований. Рейтинговая таблица на слайде 12 содержит вымышленные данные. Дополнительные слайды не входят в обязательный маршрут.
\section{Сценарий по слайдам}
\subsection{1. Введение в машинное обучение}
\timing{0–2 мин}

Представьтесь: чем занимаетесь, почему ведёте курс. Если есть собственный ML-кейс, назовите задачу и полученный результат за 60–90 секунд. Сегодня: программа и польза курса, правила оценивания, затем первая вводная по ML. Не обещайте трудоустройство после одного курса.

\subsection{2. Что вы сможете к концу семестра}
\timing{2–5 мин}

Обозначьте результат через действия студента. Например: получил таблицу, понял, что предсказывать, построил простое решение, честно оценил и объяснил его. Фокус курса — классическое ML; он создаёт базу для дальнейшего изучения нейросетей.

\subsection{3. От данных к первым моделям}
\timing{5–8 мин}

Проведите аудиторию по темам курса без календаря: сначала формулируем задачу и учимся проверять качество, затем разбираем модели и способы обучения. Математику вводим по мере необходимости. Это маршрут курса, а не список тем сегодняшней пары.

\subsection{4. От отдельных моделей к полному решению}
\timing{8–10 мин}

Продолжите маршрут: от отдельных моделей к ансамблям, затем к воспроизводимому эксперименту. Объясните связи между темами, не закрепляя их за неделями. Соревнования дадут возможность собрать изученное в самостоятельное решение.

\subsection{5. Где нужна база машинного обучения}
\timing{10–14 мин}

Назовите четыре компании. Яндекс и SberDevices публикуют темы ML-секций. У Т-Банка есть отдельная страница интервью ML-инженеров. Для Авито ссылаемся на публичную матрицу компетенций DS, а не выдаём её за перечень вопросов собеседования. У junior указаны эксперименты с метриками, признаками и гиперпараметрами. Курс даёт базу, конкретная роль может требовать дополнительные знания. Официальные источники проверены 20.10.2026.

\subsection{6. Что нужно уметь объяснять}
\timing{14–17 мин}

Это наши учебные примеры вопросов по перечисленным темам, не дословные вопросы конкретной компании. Не требуйте ответов сегодня. Пример для 99\%: если мошеннических операций 1\%, ответ «всегда обычная» даёт 99\% accuracy и не находит мошенничество. Умение объяснить базу полезно и при выборе готовых решений в работе.

\subsection{7. Базовый максимум за семестр: 100 баллов}
\timing{17–21 мин}

Проговорите каждую составляющую. Четыре опроса дают 20, два решения 50, две защиты 30. Это абсолютная сумма баллов, а не оценка относительно места одногруппника. Не обещайте дополнительные баллы за посещаемость или обычные практики: таких правил нет в документе курса. Источник: syllabus/grading\_and\_homework.md.

\subsection{8. Два соревнования и две защиты}
\timing{21–24 мин}

На каждое соревнование приходится до 40 базовых баллов вместе с защитой. Место на публичном leaderboard не заменяет воспроизводимое решение и защиту. По возможности задачи будут разными: классификация и регрессия. Публичные notebooks допустимы с указанием источника; студент должен объяснить свои изменения и их эффект. Детальная рубрика частичного начисления баллов ещё не задана; опубликуйте её до старта каждой работы.

\subsection{9. Четыре коротких опроса за курс}
\timing{24–26 мин}

Объясните правило до первого оцениваемого опроса. Попросите отключить уведомления и не переключать вкладки: потеря фокуса может фиксироваться платформой. До первого опроса преподаватель проверяет таймер и блокировку на пробной попытке. Не обещайте, что сервис запрещает системные скриншоты. Источник: syllabus/quiz\_platform.md.

\subsection{10. Как баллы переводятся в оценку}
\timing{26–29 мин}

Проверьте понимание границ: 60 — 2, 61 — 3, 76 — 4, 91 — 5. Не округляйте дробные баллы произвольно: интервалы заданы явно. На следующем слайде объясните бонусы. Источник: syllabus/grading\_and\_homework.md.

\subsection{11. Учебные медали}
\timing{29–31 мин}

Это внутренние учебные медали курса, а не официальные медали Kaggle. Их критерии пока не определены: не называйте процентили, число призовых мест или минимальный score. Также до старта соревнования нужно опубликовать рубрику начисления неполных 25 и 15 баллов. Сохраняем уже записанные бонусы и правило золота. Источник: syllabus/grading\_and\_homework.md.

\subsection{12. Все составляющие оценки видны в таблице}
\timing{31–33 мин}

Это демонстрационные данные, а не действующий рейтинг группы. Посчитайте строку Б вслух: 16+41+25+5=87, значит 4. Покажите, что студент видит, откуда берётся сумма. Для золота итоговая 5 выставляется по правилу медали независимо от суммы; в примере золотых медалей нет. Конкретная ссылка на рабочий рейтинг будет выдана отдельно, если таблица ещё не создана.

\subsection{13. Что вообще делает машинное обучение?}
\timing{33–36 мин}

Правила и модели не противоположны абсолютно: обучение — тоже программа с заранее выбранным алгоритмом. Отличие в том, что конкретная зависимость определяется данными. Спросите: где всё-таки нужен человек? Ответ: постановка, данные, метрика, решение о применении.

\subsection{14. Когда ML вообще уместен?}
\timing{36–38 мин}

После определения ML задайте вопрос «почему не просто if/else?». Если задача строго описывается коротким устойчивым правилом, нет смысла усложнять. ML полезен, когда зависимость сложна, есть данные и можно измерить выигрыш. Пример: проверка возраста по дате рождения — правило; оценка вероятности оттока — кандидат для ML. Подчеркните ограничения: сдвиг данных, приватность, стоимость ошибок.

\subsection{15. Можно ли предсказать цену квартиры?}
\timing{38–41 мин}

Начните вводную с конкретной задачи: по сведениям о квартире предсказать цену. Попросите в чате назвать недостающие данные: район, состояние, дата, тип дома. Названия типов задач введём после примера. Выбор алгоритма не заменяет хорошую постановку.

\subsection{16. Что модель получает при обучении}
\timing{41–44 мин}

Строка — объект. Площадь, этаж, район — признаки. Цена — целевая переменная. У новой квартиры цена неизвестна. Уточните, что реальный датасет требует обработки категориальных признаков, проверки пропусков и происхождения данных. Таблица только иллюстрирует постановку.

\subsection{17. Есть правильный ответ для обучения?}
\timing{44–47 мин}

Сопоставьте с квартирой: известные цены прошлых сделок — учитель. Без target можно, например, искать группы похожих квартир. Это два режима обучения, а не две библиотеки. Кратко упомяните обучение с подкреплением: агент учится выбирать действия по награде; оно за пределами основного курса.

\subsection{18. Ответ — число или категория?}
\timing{47–50 мин}

Дайте аудитории два контрастных примера и попросите определить тип. Формулы здесь лишь фиксируют тип результата: действительное число или один из конечного набора классов. Классификатор может также выдавать вероятности классов — это отдельно разберём позже. Если кто-то спрашивает про прогноз рейтинга 1–5, допустимы разные постановки; зависит от смысла и метрики.

\subsection{19. Ищем структуру без целевой переменной}
\timing{50–52 мин}

Покажите разные цели без учителя: кластеризация, снижение размерности, поиск аномалий. Кластер — результат выбранных признаков и меры близости, а не заранее известный класс. K-Means и DBSCAN группируют объекты; PCA уменьшает число координат. Подробности на неделе 15. Обнаружение аномалий также бывает с учителем, если есть размеченные примеры.

\subsection{20. Какую задачу здесь решают?}
\timing{52–55 мин}

Дайте минуту на ответы в чате до раскрытия фрагмента. Затем спросите не только название, но и признак решения: у А ответ числовой, у Б известные классы, у В меток нет. Если ответы расходятся, повторите слайд про обучение с учителем и без, а не двигайтесь дальше механически.

\subsection{21. «Классификация» — задача, «дерево» — алгоритм}
\timing{55–57 мин}

Для спама можно попробовать логистическую регрессию, дерево, бустинг, нейросеть. Название задачи не диктует модель. И наоборот: KNN может решать и регрессию, и классификацию. Это частый источник путаницы у новичков. Спросите, может ли дерево предсказывать цену квартиры; ответ да — дерево регрессии.

\subsection{22. Разные способы построить предсказание}
\timing{57–59 мин}

Короткий обзор механизмов, без вывода формул и сравнения победителей. Все эти семейства допускают варианты для регрессии и классификации. Обсудите один пример: KNN ищет похожие квартиры, дерево задаёт пороговые вопросы. Остальные методы разберём в соответствующих неделях.

\subsection{23. Как модель учится на примерах}
\timing{59–62 мин}

Объясните смысл обучения без команд библиотеки. В качестве примера возьмите линейную модель цены по площади: сначала у неё неудачный наклон, затем алгоритм находит более подходящие коэффициенты. Функция потерь задаёт, какие ошибки штрафовать. Не каждый алгоритм буквально повторяет эти шаги: линейная модель может решать систему уравнений, дерево выбирает разбиения, KNN сохраняет примеры. Снижение ошибки на обучении ещё не гарантирует качества на новых квартирах.

\subsection{24. Недообучение и переобучение}
\timing{62–65 мин}

Спросите, почему модель с нулевой ошибкой на учебных квартирах может оказаться бесполезной. При недообучении выбранная модель не улавливает существенную зависимость. При переобучении она использует случайные особенности обучающей выборки. Пример с одной ценой иллюстрирует слишком грубое решение для неоднородных квартир, а не утверждает, что константа всегда плоха. Разрыв ошибок также бывает из-за сдвига распределения. Следующий слайд объясняет механизмы через bias и variance.

\subsection{25. Bias–Variance Tradeoff}
\timing{65–72 мин}

Предложите мысленный эксперимент: много раз собираем разные обучающие наборы одинакового размера из одного распределения и предсказываем цену одной и той же квартиры. Bias — отличие среднего прогноза от истинной средней цены при этих признаках. Variance — разброс прогнозов между такими обучениями. Грубая модель может стабильно ошибаться, гибкая — сильно реагировать на состав train. На схеме Bias\textsuperscript{2} падает, variance растёт, их сумма с неустранимым шумом даёт ожидаемую квадратичную ошибку. Шум показан пунктирной горизонталью. Это иллюстрация типичного режима, не универсальный закон формы кривых. Не называйте variance разбросом цен квартир или ошибками одной модели. При переобучении пробуем ограничение глубины, регуляризацию или больше независимых данных; при недообучении — более выразительную модель и полезные признаки. Выбор проверяем на validation. Источники: https://www.teach.cs.toronto.edu/~csc311notes/bias\_variance\_decomposition/bv.html и https://scikit-learn.org/stable/auto\_examples/ensemble/plot\_bias\_variance.html .

\subsection{26. Успех на знакомых квартирах ещё не успех}
\timing{72–75 мин}

Аналогия: проверять студента по ответам, которые он только что заучил, бессмысленно для оценки переноса знаний. В реальной практике схемы разделения зависят от данных, времени и объёма выборки; здесь только принцип. Нельзя подбирать гиперпараметры по test. Позже вернёмся к этому подробно.

\subsection{27. Утечка данных · Data Leakage}
\timing{75–79 мин}

Покажите четыре механизма утечки. Комиссия, равная проценту итоговой цены, косвенно раскрывает target. Действия клиента после даты прогноза раскрывают будущее. Копии одного объявления в разных выборках позволяют узнать учебный объект. Средние для заполнения пропусков и отбор признаков нельзя вычислять по всей выборке до split: validation и test начинают влиять на обучение. Если предсказываем для новых клиентов, делим по клиентам; если будущее — по времени. Связанные объекты опасны, когда такое пересечение не соответствует сценарию применения. Постоянный подбор модели по test тоже загрязняет оценку. Высокое качество само по себе не доказывает утечку.

\subsection{28. Модель должна быть полезнее простого прогноза}
\timing{79–82 мин}

Медиану вычисляем только по train. Для MAE медиана — естественный константный baseline: она минимизирует сумму абсолютных отклонений на train. Попросите посчитать (1+2+1)/3 и назвать единицу измерения. Числа учебные. Модели сравниваем по validation, а test оставляем для финальной оценки выбранного решения. Средняя цена тоже возможный baseline; она оптимальна среди констант для квадратичной ошибки, не для MAE.

\subsection{29. Параметр ≠ гиперпараметр}
\timing{82–84 мин}

Не погружайтесь в вывод линейной модели. Верхнее T — запись скалярного произведения. Спросите, что из примеров будет искать алгоритм, а что нужно выбрать до fit. В строгом смысле выбор гиперпараметров сам является процедурой настройки по данным, но это не внутренние параметры одной обученной модели. Test для такого выбора не используем.

\subsection{30. Полный маршрут ML-решения}
\timing{84–86 мин}

Пройдите маршрут на квартире. Метрику и способ проверки определяем при постановке, до выбора модели; положение слова «метрика» в цепочке означает вычисление результата. Данные сначала разделяем, статистики preprocessing и baseline считаем на train. Повторное улучшение опирается на validation, test используем в конце. Baseline в сегодняшнем примере — медиана train.

\subsection{31. Что разобрали сегодня}
\timing{86–90 мин}

Кратко свяжите разобранное в одну картину: сформулировали задачу, нашли признаки и target, выбрали тип задачи и обсудили обучение. Качество проверяем на новых данных, сравниваем с baseline и следим за утечками. Сложность выбираем с учётом недообучения и переобучения. Завершите ответами на вопросы аудитории.

\newpage
\appendix
\section{Дополнительные объяснения для подготовки}
Ниже сохранены подробные объяснения и математические углубления. Их порядок тематический, а не порядок показа. Основной сценарий и тайминг заданы выше.
\section{Что такое ML и когда оно вообще нужно}



\term{Обычная программа} содержит правила, определённые разработчиком. \term{Обучение по данным} выбирает конкретную зависимость из семейства допустимых моделей, используя наблюдаемые примеры. При этом ML не заменяет постановку задачи: человек по-прежнему выбирает, что считать объектом, какие данные использовать, какую ошибку считать критичной и какое действие предпринять после прогноза.

Сравнение, которое можно произнести: вычислить возраст по дате рождения --- простое правило; предсказать вероятность ухода клиента по истории поведения --- кандидат для ML. Не противопоставляйте ML и правила абсолютно: модель сама является программой, а признаки и пороги часто включают ручные решения.

\begin{saybox}{Вопрос аудитории}
«Если формула стоимости доставки заранее известна и стабильно работает, зачем нам обучать модель?» Ожидаемый ответ: незачем, если правило покрывает задачу и удовлетворяет требованиям. ML добавляет сложность, стоимость поддержки и риск ошибок.
\end{saybox}

Перед выбором алгоритма задайте три вопроса. \textbf{Есть ли полезное решение после прогноза?} \textbf{Есть ли репрезентативные, допустимые к использованию данные?} \textbf{Есть ли способ сравнить результат с простым решением?} Если один ответ отрицательный, «попробуем бустинг» не является полноценным планом.

\section{С учителем и без: две разные постановки}



\subsection{Обучение с учителем}

Для каждого учебного объекта есть пара $(x_i,y_i)$: описание $x_i$ и известный ответ $y_i$. Цель --- построить правило $f$, которое на новом $x$ выдаёт прогноз $\hat y=f(x)$. «Учитель» здесь означает метки в данных, а не человека, который стоит над алгоритмом.

Идеальная математическая цель:
\[
f^*=\arg\min_{f\in\mathcal F}\;\mathbb E_{(x,y)\sim P}\bigl[L(f(x),y)\bigr].
\]
Здесь $\mathcal F$ --- семейство кандидатов; $L$ --- цена ошибки на одном объекте; $P$ --- неизвестное распределение будущих случаев. Объясните словами: мы хотим, чтобы модель ошибалась мало \emph{на будущих объектах}, а не только на тех, которые уже лежат в таблице. Распределение $P$ неизвестно, поэтому используем конечную выборку как его приближение. Если будущее существенно изменится, гарантия переноса становится слабее.

\term{Регрессия}: ответ --- численная величина, например цена квартиры или время доставки. Обычно пишут $\hat y\in\mathbb R$. \term{Классификация}: ответ принадлежит конечному набору классов, например спам/не спам, $\hat y\in\{0,\ldots,K-1\}$. Классификатор может выдавать также вероятности классов; окончательный класс получается после правила принятия решения, часто с порогом. Прогноз рейтинга 1--5 может допускать несколько постановок --- важны смысл шкалы и метрика.

\begin{warnbox}{Не путать задачу и алгоритм}
«Классификация» и «регрессия» отвечают на вопрос, \emph{какой результат нужен}. «KNN», «линейная модель» и «дерево» отвечают на вопрос, \emph{как строить предсказание}. Дерево может быть классификатором или регрессором; KNN --- тоже.
\end{warnbox}

\subsection{Обучение без учителя}

У объектов есть признаки $x_i$, но нет заданного целевого ответа $y_i$. Алгоритм ищет структуру: \term{кластеризация} группирует похожие объекты, \term{снижение размерности} даёт более компактное представление, \term{обнаружение аномалий} ищет необычные случаи. K-Means и DBSCAN --- два алгоритма кластеризации; PCA --- метод снижения размерности. Не представляйте PCA как разновидность классификации.

Кластер не тождествен классу: он возникает из выбора признаков, масштаба и понятия близости. Разбив клиентов по сумме покупок и частоте посещений, можно получить одни группы; по категориям товаров --- другие. Название «лояльные клиенты» даёт аналитик после исследования, а не алгоритм как объективную истину.

\begin{saybox}{Проверка понимания перед следующей частью}
«Длительность поездки по размеченной истории --- что это? Спам по размеченным письмам? Группы клиентов без меток?» Ожидаемые ответы: регрессия; классификация; кластеризация. Попросите объяснить признак выбора, а не просто назвать термин.
\end{saybox}

\section{Карта алгоритмов без ложной гонки за «лучшим»}



У каждого семейства есть своё полезное предположение о данных. \term{KNN}: близкие в выбранной метрике объекты часто имеют похожие ответы; результат зависит от масштабирования и выбора $k$. \term{Линейные модели}: прогноз выражается через взвешенную сумму признаков; сложные зависимости требуют подходящего представления признаков. \term{Деревья}: последовательные пороговые вопросы разбивают пространство; одиночное глубокое дерево может быть нестабильным и переобучаться. \term{Случайный лес}: усредняет множество различающихся деревьев, снижая их нестабильность. \term{Бустинг}: строит ансамбль последовательно, добавляя модели, которые исправляют ошибки текущего решения. \term{SVM}: ищет разделяющую границу с запасом; детали будут позже.

Нет универсально лучшего алгоритма. Уместность зависит от объёма и природы данных, требований к интерпретации, скорости предсказания, стоимости ошибок и результатов честной проверки. На первой лекции достаточно показать принципы, а не выводить каждый метод.

\begin{saybox}{Как сказать про \texttt{fit()}}
«Интерфейс библиотеки выглядит одинаково: \texttt{fit(X,y)}. Но за ним скрыты разные действия. KNN в основном сохраняет примеры; линейная модель оценивает коэффициенты; дерево выбирает пороги разбиений. Команда одна, математическая процедура разная».
\end{saybox}

\section{Сквозной пример: цена квартиры}



В основной презентации конкретный пример вводится до карты задач. При повторном обсуждении спросите: \emph{«Цена квартиры --- это какой тип задачи?»} Ответ: регрессия с учителем, если есть история сделок с известными ценами. Затем покажите объект: квартира площадью 62 м$^2$, 7-й этаж, 12 минут до метро. Попросите назвать недостающие сведения: район, состояние, год постройки, дата сделки, тип дома. Три признака намеренно недостаточны.

В таблице строка --- \term{объект}, столбцы характеристик --- \term{признаки}, известная цена --- \term{целевая переменная} или target. Запись
\[
\mathcal D=\{(x_i,y_i)\}_{i=1}^{n}
\]
означает $n$ пар: $x_i$ --- вектор характеристик $i$-й квартиры, $y_i$ --- её цена. Индекс $i$ только нумерует примеры. Для новой квартиры есть $x_\mathrm{new}$, но её $y_\mathrm{new}$ пока неизвестен. В презентации три сделки показаны только для объяснения --- этого объёма недостаточно для настоящего обучения.

\subsection{Признак и утечка}

Главный тест для любого признака: \emph{«Будет ли он доступен в момент реального прогноза?»} Площадь и этаж обычно известны; финальная цена сделки --- ответ, а не признак. Информация, появившаяся после продажи, тоже недопустима для прогноза до продажи. Если она попадёт в $X$, качество в ноутбуке может быть фантастическим, а применение модели --- невозможным. Это \term{утечка данных}. Отдельная опасность: один и тот же объект или почти дубликаты оказываются и в обучении, и в проверке.

\subsection{Эмпирическая ошибка}

Поскольку неизвестное распределение будущего $P$ нам недоступно, многие методы ориентируются на ошибку в доступных данных:
\[
\hat R(f)=\frac{1}{n}\sum_{i=1}^{n} L\bigl(f(x_i),y_i\bigr).
\]
Это средняя потеря на $n$ обучающих примерах. Формула описывает общую идею оптимизации, но \textbf{не утверждает}, что каждый алгоритм буквально обучается градиентным спуском: KNN устроен иначе. Минимум $\hat R$ --- не конец истории. Можно добиться почти нулевой ошибки, запомнив шум, и плохо работать на новых квартирах.

\section{Обобщение, проверка и выбор метрики}



\subsection{Обучение и предсказание}

При обучении алгоритм видит признаки и ответы: \texttt{model.fit(X\_train, y\_train)}. При применении он получает признаки нового объекта: \texttt{model.predict(X\_new)}. Общее обозначение прогноза:
\[
\hat y=f(x;\theta).
\]
Шляпка над $y$ означает прогноз, а не истинный ответ. $\theta$ обозначает обученное состояние модели в широком смысле; это не обязательно числовой вектор весов для каждого алгоритма. У линейной модели частный случай: $\hat y=w^{\mathsf T}x+b$, где $w$ и $b$ находятся из данных. Число соседей $k$ в KNN --- гиперпараметр: его задают до конкретного обучения и выбирают по корректной проверке.

\subsection{Почему train недостаточно}

Цель --- обобщающая способность: работа на новых объектах. Проверка на тех же строках, которые модель видела при fit, измеряет главным образом способность объяснить знакомые данные. Низкая ошибка на train и высокая на test может говорить о переобучении, утечке при проверке или сдвиге данных; она не доказывает единственную причину. Высокая ошибка и на train, и на test может указывать на слишком слабую модель, плохие признаки или неправильную постановку.

Разделение данных служит разным ролям: \term{train} --- обучение и построение baseline; \term{validation} --- выбор модели, признаков и гиперпараметров; \term{test} --- заключительная оценка уже выбранного решения. Если многократно смотреть на test и менять модель, он фактически превращается в ещё одну validation-выборку.

Честный split должен имитировать реальное применение. Для будущих сделок часто важен временной порядок: обучение на прошлом, проверка на более позднем. Для повторяющихся клиентов, домов или устройств требуется учитывать группы, чтобы одни и те же сущности не подсказали ответ с обеих сторон split. Случайное разбиение не универсально.

\subsection{Baseline и метрика}

Константный baseline для цены квартиры при MAE --- медиана цен \emph{обучающих} сделок:
\[
\hat y_{\mathrm{base}}=\operatorname{median}(y_1,\ldots,y_n).
\]
Его применяют к отложенным объектам и сравнивают с более сложной моделью на одной и той же метрике. Для MAE на $m$ проверочных объектах:
\[
\operatorname{MAE}=\frac{1}{m}\sum_{j=1}^{m}\left|y_j-\hat y_j\right|.
\]
MAE выражается в тех же единицах, что и цена. Например, ошибки $1$, $2$ и $3$ млн рублей дают MAE $2$ млн рублей. Этот числовой пример полезен, потому что студенты сразу видят смысл усреднения. При MAE медиана минимизирует сумму абсолютных отклонений на train; при MSE среди констант оптимально среднее. Если сложная модель хуже baseline, задайте вопрос не «какой ещё алгоритм добавить?», а «не ошиблись ли мы с данными, признаками, разделением или метрикой?».

\term{Функция потерь} даёт сигнал для настройки параметров; \term{метрика} сообщает, насколько результат полезен по выбранному критерию. Это различие ролей, а не требование двух разных формул. Для задачи мошенничества простая accuracy может быть обманчива: при редких положительных случаях модель «всегда нет» имеет высокую долю правильных ответов, но не решает задачу. Здесь пока не выводите precision и recall --- лишь укажите, что цена ложной тревоги и пропущенного случая разная.

\begin{saybox}{Мини-разбор у доски}
Пусть истинные цены трёх квартир: 8, 10, 12 млн, а прогнозы: 9, 8, 13 млн. Абсолютные ошибки: 1, 2, 1 млн. MAE $=(1+2+1)/3\approx1{,}33$ млн. Попросите студентов самостоятельно назвать единицу измерения и подумать, как оценить, достаточно ли это хорошо: нужен контекст рынка и baseline.
\end{saybox}

\section{Заключение и первая практика}



Пройдите полную цепочку ещё раз, но теперь студенты должны объяснять каждый шаг: \textbf{задача} (цена), \textbf{данные} (история сделок), \textbf{разделение} (имитация будущего), \textbf{baseline} (медиана цен train), \textbf{модель} (например, линейная), \textbf{метрика} (MAE), \textbf{анализ ошибок} (где промахиваемся), \textbf{вывод} (что изменить и почему). Цепочка итеративна: обнаружив проблемы, мы можем вернуться к задаче и данным.

На первой практике студенты в Jupyter загрузят небольшой датасет, выделят $X$ и $y$, разделят данные, обучат baseline, получат предсказания и посчитают метрику. До неё им следует проверить Python, Jupyter, Git и доступ к GitHub. При проблемах с локальной средой Colab можно временно использовать как резерв.

\begin{saybox}{Три финальных вопроса аудитории}
1. Цена квартиры --- задача с учителем или без, регрессия или классификация?\\
2. Группы клиентов без меток --- чем отличаются от заранее заданных классов?\\
3. Почему высокое качество на train не доказывает, что модель хорошо работает?
\end{saybox}

\section{Шпаргалка преподавателя: частые ошибки и точные ответы}

{\small\begin{longtable}{@{}P{52mm}P{105mm}@{}}
\toprule
Неверное утверждение / вопрос & Что ответить\\
\midrule
\endhead
«ML сам найдёт правильную задачу» & Нет: target, доступные признаки, критерий качества и действие после прогноза задаются отдельно.\\
«Классификация --- это алгоритм» & Это тип задачи. Её можно решать разными моделями.\\
«Кластер --- это настоящий класс» & Кластер зависит от выбранных признаков, масштаба, метрики и метода; смысл требует отдельной интерпретации.\\
«\texttt{fit()} всегда меняет веса» & Нет: у KNN он в основном сохраняет примеры; у линейной модели оценивает коэффициенты; дерево строит разбиения.\\
«0 ошибок на train --- идеальная модель» & Это может быть переобучение или утечка. Проверяем на данных, не участвовавших в обучении и подборе.\\
«Можно подобрать модель по test» & Тогда test перестаёт быть независимой заключительной проверкой. Выбор делаем по validation или CV.\\
«Baseline должен быть сложным» & Напротив: это понятная простая точка сравнения, которую трудно ошибочно улучшить только на бумаге.\\
«Метрика всегда равна loss» & Иногда формулы совпадают, иногда нет; различайте роль при оптимизации и роль при оценке полезности.\\
«Какой алгоритм лучший?» & Без постановки, данных, метрики и ограничений такой вопрос не имеет содержательного ответа.\\
\bottomrule
\end{longtable}}

\section{Источники}
\begin{itemize}
\item Программа, оценивание и опросы: документы в папке \texttt{syllabus} репозитория курса.
\item \href{https://yandex.ru/jobs/interview/mldev?turbo=true}{Яндекс: отбор ML-специалистов}.
\item \href{https://sberdevices.ru/career/ML/interview/}{SberDevices: технические секции}.
\item \href{https://www.youtube.com/watch?v=k3UJOG-DKHE}{Яндекс: Тренировки по ML 1.0, вводная лекция} как методический ориентир.
\end{itemize}
Проверка внешних источников: 20 октября 2026. Примеры вопросов в основной презентации учебные и не являются дословными вопросами компаний.
\end{document}
